\section{Overview}
\label{sec:overview}
Before presenting the full details of our proposed type system, we
begin with a tour of its key features and how they enable both safety
and reachability reasoning.

%% Readers can see this
%% connection via the type erasure function $\eraserf{\tau}$, which
%% converts context types to basic types by erasing these refinements. In
%% the following paragraphs, we introduce these refinements in detail via
%% examples.

\paragraph{Demonic and Angelic Modalities}
% BD: The goal of this subsection is to equip the reader with the
% right intuitions of what the modalities mean when ascribed to types
% and when used in the typing context of a typing judgment. I'm
% currently of the opinion that the lower / upper bound framing is the
% right intuition (particularly for the type context case), as opposed
% to the angelic / demonic one.

% The reader should have also internalized the motto'that "angelic
% refinement = reachability, demonic refinement = safety" by the end
% of this subsection.

% Note that fundamental question that our type system answers is: what
% does the image of this term look like, when it is evaluated under a
% particular restriction? One of the interesting things about the type
% system is a) the sorts of restrictions our type contexts can
% describe, and b) how those restrictions afford the type certain
% reasoning capabilities when verifying safety / reachability
% properties.

As in other refinement type systems, a context type refines a base
type $b$ with a qualifier $\phi$, a sentence in first-order
logic. Every context type is equipped with a modality indicating
whether it should be interpreted either demonically ($\ort{b}{\phi}$)
or angelically ($\urt{b}{\phi}$). When ascribed to a term $e$, a
context type describes both safety and reachability properties of
$e$. As discussed in \autoref{sec:intro}, the demonic context type
$\ort{b}{\phi}$ allows $e$ to (at most) acquire any value $v$ that
satisfies its qualifier $\phi$.  Conversely, an angelic context type
$\urt{b}{\phi}$ stipulates that $e$ must (at least) acquire every
value satisfying $\phi$.

% Informally, we have
% \begin{align*}
%   &\vdash e : \ort{b}{\phi} \quad\text{iff}\quad \forall v. \mstep{e}{v} \implies \phi[\vnu \mapsto v]
%   \\&\vdash e : \urt{b}{\phi} \quad\text{iff}\quad \forall v. \phi[\vnu \mapsto v] \implies  \mstep{e}{v}
% \end{align*}\noin

To illustrate this difference, consider the following typing
judgements which ascribe different context types to the constant term
$1$:
\begin{align*}
  &\vdash 1 : \ort{\Nat}{\top}\;\judgok
  && \vdash 1 : \urt{\Nat}{\top}\;\judgfail
  \\&\vdash 1 : \ort{\Nat}{\vnu > 0}\;\judgok
  && \vdash 1 : \urt{\Nat}{\vnu > 0}\;\judgfail
  \\&\vdash 1 : \ort{\Nat}{\vnu = 1}\;\judgok
  && \vdash 1 : \urt{\Nat}{\vnu = 1}\;\judgok
\end{align*}\noindent
The judgements on the left assert \emph{safety} properties, similar to
a standard refinement type systems, where the constant $1$ needs to
satisfy the corresponding qualifier. The judgements on the right, in
contrast, describe \emph{reachability} constraints: since the constant
$1$ cannot reach every natural number or positive number, the first
two judgements are invalid. The last (singleton) type precisely
captures the reachability properties the term exhibits. Under an empty
typing context, a deterministic term\footnote{All terms in \langname{}
  are deterministic, a restriction that we will lift in
  \autoref{sec:extension}.} can only be assigned either an empty or
singleton angelic type, as it can evaluate to at most one
value. Richer reachability properties naturally arise when reasoning
about the behavior of a term within a broader program context. To do
so, our type system \emph{holistically} considers the \emph{set} of
values that a term can evaluate to, i.e., its image, under a
restricted \emph{set} of environments, or \emph{world}, that are
constrained by a typing context.
% Put another way, our type system describes the image of a term under
% a set of environments consistent with a typing context.

Intuitively, a demonic qualifier places an upper bound on the image of
a term, while an angelic qualifier describes a lower bound. Similar
intuitions apply to the restrictions that a typing context imposes on
a set of input environments in our type system. To see how,%  these
% bounds manifest in the guarantees provided by our type system,
consider the following typing judgements involving non-singleton
reachability properties: %
\begin{align*}
  x{:}\urt{\Nat}{x > 3} \vdash x - 3 : \urt{\Nat}{\vnu > 0}\;\judgok \\
  x{:}\urt{\Nat}{x > 4} \vdash x - 3 : \urt{\Nat}{\vnu > 0}\;\judgfail
\end{align*} %
\noindent
Clearly, for any particular positive number, the term $x - 3$ will
evaluate to that number, assuming an appropriate assignment to
$x$. The angelic types of $x$ in each typing context place different
restrictions on the set of possible assignments to $x$ in a world:
under the first, for every $n$ greater than $3$, the world must
contain an environment in which $x$ is $n$; the second allows a world
which does not include a mapping of $x$ to $4$, and thus cannot
guarantee the image of $x - 3$ includes $1$.

Switching the modality on the type of $x$ inverts the direction of
these bounds, thereby restricting the set of assignments to $x$ that
are impossible. As a consequence, the reachability guarantees that our
type system can provide about deterministic terms involving only
demonically qualified variables are somewhat limited:
\begin{align*}
  x{:}\ort{\Nat}{x > 2} \vdash x - 3 : \urt{\Nat}{\vnu > 0}\;\judgfail \\
  x{:}\ort{\Nat}{x > 4} \vdash x - 3 : \urt{\Nat}{\vnu > 0}\;\judgfail
\end{align*} %
Observe there are many singleton worlds containing in which $x$ is
assigned a value greater than $2$, but $x - 3$ can only reach a single
result value under such a world!
% If all we know about $x$ is that it cannot be less than $2$, after
% all, there's little we can say about the values that $x-3$ can reach!
Similarly, our type system can provide somewhat limited safety
guarantees when only equipped with reachability assumptions:
\begin{align*}
  x{:}\urt{\Nat}{x > 4} \vdash x - 3 : \ort{\Nat}{\vnu > 0}\;\judgfail \\
\end{align*} %
The image of $x - 3$ includes a non-positive value in any world
containing an environment where $x$ is $4$, but the lower bound
provided by the angelic type of $x$ in the typing judgement above does
not preclude such an environment. Using our type system, it is still
possible in to derive meaningful safety guarantees using reachability
assumptions through the use of $\oplus$ operator, as we shall see
shortly.

\BD{At some point early on, we need to introduce the notion of method
  predicates. This may just happen organically with the addition of
  the buggy merge example to the introduction... }

% This
% A typing
% judgement can ascribe a non-trivial (i.e., non-singleton) angelic type
% to its term based on its typing context.  For example, the following
% type judgement is valid:
% \begin{align*}
%   x{:}\urt{\Nat}{\top} \vdash x - 3 : \urt{\Nat}{\vnu > 0}\;\judgok
% \end{align*}\noindent
% Clearly, the term $x - 3$, when evaluated, can produce all positive
% numbers because its corresponding type context $x{:}\urt{\Nat}{\top}$
% assumes that assignments to variable $x$ can reach all natural
% numbers; in particular, this means $x$ can be any number greater than
% 3.  An angelic type context allows the type system to presume a
% desirable assignment of $x$ depending on the shape of the
% right-hand-side type; e.g., in order to reach $n$, the assignment of
% $x$ should be $n + 3$.

% We assume an \emph{external impure environment} that serves as a source
% of angelic inputs.  For example, a test generator that yields random
% numbers can serve as an example of such a source:
% \begin{align*}
%   \Envir_x(e) \doteq \zlet{x}{\Code{nat\_gen()}}{\boxed{e}}
% \end{align*}
% For our purposes here, we do not include environments such as
% $\Envir_x$ as part of our core language.

\paragraph{Subtyping}
% \BD{TBH, the material in this subsection does not seem
%   particularly novel (the subtyping relations are similar to prior
%   work) or important to what follows (I don't believe swapping
%   modalities on singleton refinements is not used elsewhere). }

The subtyping relationship between two angelic or two demonic
refinements is similar to contemporaneous refinement type
systems~\cite{LiquidHaskell, CoverageType}, where the subtype
relationship between demonic refinements is covariant with the
qualifier, and that of angelic refinements is contravariant with the
qualifier. These relationships align with the interpretation of
angelic and demonic modalities as lower- and upper- bounds: a term
which can reach all positive numbers can also reach all numbers
greater than 1, while a term that always safely evaluates to a number
greater than 1 is also guaranteed to evaluate to a positive number.
\begin{align*}
  &\ort{\Nat}{\vnu > 0} \subty \ort{\Nat}{\vnu > 1}\;\judgok
  &&\urt{\Nat}{\vnu > 1} \subty \urt{\Nat}{\vnu > 0} \;\judgfail \\
  & \ort{\Nat}{\vnu > 0} \subty \ort{\Nat}{\vnu > 1}  \;\judgfail
  && \urt{\Nat}{\vnu > 0} \subty \urt{\Nat}{\vnu > 1}\;\judgok
\end{align*} %
\noindent These ideas extend naturally to a refinement relation
between typing contexts, $\Gamma_1 \ctsubseteqA \Gamma_2$, which
intuitively reads as ``any term that has type $\tau$ under $\Gamma_1$
can also be typed with $\tau$ under $\Gamma_2$''.

The subtyping relationship between angelic and demonic refinement
types is more nuanced. In general, no subtyping relationship exists
between two arbitrary refinements with different modalities:
\begin{align*}
  &\ort{\Nat}{\vnu > 0} \subty \urt{\Nat}{\vnu > 1}\;\judgfail
  & &\urt{\Nat}{\vnu > 0} \subty \ort{\Nat}{\vnu > 0} \;\judgfail
\end{align*}\noindent
Knowing that a term can only evaluate to positive numbers does not
allow us to draw any conclusions about which positive numbers it
reaches, for example. Similarly, simply knowing that a term can reach
every positive number does not preclude it from also evaluating to 0.
\begin{wrapfigure}{r}{0.48\linewidth}
  \centering
  \begin{minted}[linenos = false, fontsize = \footnotesize, escapeinside=??]{OCaml}
val g : ?$\Code{i}{:}\ort{\Int}{\vnu \le 10}\ctsarr\ort{\Int}{\vnu \neq 0}$?

?$\vdash$? let f (x: int) (y: int) =
  if y > 10 then (g x) + 1
    else x + 1
: ?$\Code{x}{:}\urt{\Int}{x \le 5}\ctsarr\Code{y}{:}\urt{\Int}{y \le 15}\ctsarr\urt{\Int}{\vnu = 1}$?
  \end{minted}
  \caption{A simple example of how may-analysis helps
    must-analysis~\cite{CompositionalMayMust}.}
  \label{fig:may-must-reach}
\end{wrapfigure}
That said, numerous prior works~\cite{...} have shown that knowledge
about the safety properties of terms can aid the automated
verification of reachability properties, and vice-versa. When
establishing that the function \Code{f} in
\autoref{fig:may-must-reach} must return \Code{1} on some arguments,
for example, the program analysis of \citet{CompositionalMayMust} uses
a not-may summary analogous to the signature of \Code{g} to conclude
calling it will never produce the \Code{0} value needed for \Code{f}
to yield \Code{1}. As we shall see shortly, our type system also
naturally supports using angelic information in the type context to
verify a meaningful form of qualified safety properties, similar to
that found in \citet{ZDS23} and \citet{RHLE}.

% It is only possible to switch the modality of a refinement when its
% qualifier is a singleton, or when the supertype qualifier is trivial:
% \begin{align*}
%   &\ort{\Nat}{\vnu = 3} \subty \urt{\Nat}{\vnu =
%     3}\;\judgok
%   & & \urt{\Nat}{\vnu = 3} \subty \ort{\Nat}{\vnu =
%       3}\;\judgfail \\
%   &x{:}\urt{\Nat}{\vnu > 0} \subty x{:}\ort{\Nat}{\top}\;\judgok
%   & &x{:}\ort{\Nat}{\vnu > 0} \ctsubseteqA x{:}\urt{\Nat}{\bot}\;\judgok
% \end{align*}\noindent
% When the qualifier is a singleton, safety and reachability coincide
% --- there is only one possible value, so the demonic constraint (the
% value must be 3) and the angelic constraint (the value 3 must be
% reachable) make identical demands on the term. Converting to a
% supertype with a trivial qualifier, in contrast, simply throws away
% the bounds on a term entirely.

% that an angelic context requiring $x$ to reach all numbers greater
% than 1 is a subtype of one requiring $x$ to reach all positive
% numbers. This holds because any term that is typed under the
% assumption that $x$ can reach every number greater than 1 can
% certainly be given the same type under the assumption that $x$ can
% reach every positive number, since the latter subsumes the former.
% In other words, a narrower angelic qualifier (fewer reachable
% values) is a weaker assumption, while a wider angelic qualifier
% (more reachable values) is stronger — the opposite of the demonic
% case.

% Having both angelic as well as demonic modalities in a type context
% requires a more refined subtype relation than usual.  Read
% ``$\Gamma_1 \ctsubseteqA \Gamma_2$'' as ``any term that has type
% $\tau$ under $\Gamma_1$ can also be typed with $\tau$ under $\Gamma_2$
% .''  Then, the following statements hold:
% \begin{align*}
%   &x{:}\ort{\Nat}{\top} \ctsubseteqA x{:}\ort{\Nat}{\vnu > 0}\;\judgok
%   &&x{:}\urt{\Nat}{\top} \ctsubseteqA x{:}\urt{\Nat}{\vnu > 0} \;\judgfail
%   \\&x{:}\ort{\Nat}{\vnu > 1} \ctsubseteqA  x{:}\ort{\Nat}{\vnu > 0}  \;\judgfail
%   &&x{:}\urt{\Nat}{\vnu > 1} \ctsubseteqA x{:}\urt{\Nat}{\vnu > 0}\;\judgok
% \end{align*}\noindent
% A demonic (over-approximate) context has a subtype relation that is
% covariant with the qualifier, while the angelic (under-approximate)
% context is contravariant with the qualifier.  For example, the last
% statement says that an angelic context requiring $x$ to reach all
% numbers greater than 1 is a subtype of one requiring $x$ to reach all
% positive numbers. This holds because any term that is typed under the
% assumption that $x$ can reach every number greater than 1 can
% certainly be given the same type under the assumption that $x$ can
% reach every positive number, since the latter subsumes the former.  In
% other words, a narrower angelic qualifier (fewer reachable values) is
% a weaker assumption, while a wider angelic qualifier (more reachable
% values) is stronger — the opposite of the demonic case.

% Additionally, the subtype relation must allow context shifts between
% these two modalities since they serve as resources used by the type
% system.  Specifically, a demonic context cannot be treated as a
% non-trivial angelic context:
% \begin{align*}
%   &x{:}\ort{\Nat}{\vnu > 0} \ctsubseteqA x{:}\urt{\Nat}{\top}\;\judgfail
%   &&x{:}\ort{\Nat}{\vnu > 0} \ctsubseteqA x{:}\urt{\Nat}{\vnu > 1}\;\judgfail
%   \\&x{:}\ort{\Nat}{\vnu = 3} \ctsubseteqA x{:}\urt{\Nat}{\vnu = 3}\;\judgok
%   &&x{:}\ort{\Nat}{\vnu > 0} \ctsubseteqA x{:}\urt{\Nat}{\bot}\;\judgok
% \end{align*}\noindent
% For example, the demonic context $x{:}\ort{\Nat}{\vnu > 0}$ only
% constrains $x$ to be \emph{some} positive number, providing no
% guarantee about which positive numbers are reachable.  Treating it as
% a non-trivial angelic context would require it to serve as a witness
% that \emph{every} value satisfying the angelic qualifier is reachable,
% a guarantee that a demonic context is fundamentally incapable of
% providing.   Demonic and angelic contexts align when
% their type qualifier is trivially false ($\bot$) or is a singleton
% (e.g., $\vnu = 3$). When the qualifier is a singleton, safety and
% reachability coincide — there is only one possible value, so the
% demonic constraint (the value must be 3) and the angelic constraint
% (the value 3 must be reachable) make identical demands on the term.
% Finally, an angelic context can be viewed as a subtype
% of a demonic one provided every value reachable under the angelic
% qualifier also satisfies the demonic qualifier.
% \begin{align*}
%   x{:}\urt{\Nat}{\vnu > 0} \ctsubseteqA x{:}\ort{\Nat}{\top}\;\judgok
%   &&x{:}\urt{\Nat}{\vnu > 0} \ctsubseteqA x{:}\ort{\Nat}{\vnu > 1}\;\judgfail
% \end{align*}\noindent
% The statement on the right is incorrect because the angelic context
% guarantees that $x$ can reach every positive number, including
% $1$. However, the demonic context $x{:}\ort{\Nat}{\vnu > 1}$ only
% permits values strictly greater than 1 so a term typed as the
% singleton $1$ under the angelic context would violate the demonic
% constraint.

\paragraph{Type Context Composition}
% Here we want to transition from thinking of type contexts as
% constraining the surrounding environment to thinking of them as
% affording us certain reasoning capabilities. Importantly, if we want
% to reason compositionally, we need to have a finer-grained notion of
% the sorts of capabilities it can provide
When they appear in a typing context, the bounds imposed by our two
modalities afford us different reasoning capabilities: when
establishing the reachability properties of a term, for example, an
angelic refinement lets us \emph{choose} a favorable assignment for a
variable in that term, while a demonic refinement allows us to
\emph{rule out} certain assignments to that variable showing that a
term is safe.
% BD: If there's room, we could point back to specific examples in the
% previous subsection to emphasize this point.
Importantly, the bounds imposed by the bindings of individual
variables in a typing context can interact with each other, altering
the capabilities afforded by the typing context as a whole. To
illustrate, consider the following pair of program contexts:
\begin{align}
  &\zlet{x}{\Code{nat\_gen () }}{\zlet{y}{x + 1}{\square}}
    \tag{C$_{en}$}\label{ex:EnCon} \\
  &\zlet{x}{\Code{arbitrary}}{\zlet{y}{\Code{nat\_gen()}}{\square}}
    \tag{C$_{dis}$}\label{ex:DisjCon}
\end{align}\noindent
Both contexts use \Code{nat\_gen ()} to (nondeterministically) generate
an arbitrary value for $x$. The first context then sets $y$ based on
the value as $x$, while the second opts to generate a fresh value for
$y$ instead. The four programs that result from plugging $x$ and $y$
into these contexts can all reach reach every possible positive
number, as can the program \ref{ex:DisjCon}[$x - y$].  The same can
not be said of \ref{ex:EnCon}[$x - y$], however, because the values of
$x$ and $y$ in this program are somehow \emph{entangled}.
\begin{align*}
  &\emptyset \vdash \text{\ref{ex:DisjCon}}[x] : \urt{\Nat}{\vnu > 0}
    \;\judgok \quad
  &\emptyset \vdash \text{\ref{ex:EnCon}}[x] : \urt{\Nat}{\vnu > 0}
    \;\judgok  \\
  & \emptyset \vdash \text{\ref{ex:DisjCon}}[y] : \urt{\Nat}{\vnu > 0}
  \;\judgok \quad
  &\emptyset \vdash \text{\ref{ex:EnCon}}[y] : \urt{\Nat}{\vnu > 0}
    \;\judgok \\
  & \emptyset \vdash \text{\ref{ex:DisjCon}}[x - y] : \urt{\Nat}{\vnu
    > 0} \;\judgok \quad & \emptyset \vdash \text{\ref{ex:EnCon}}[x - y] : \urt{\Nat}{\vnu > 0} \;\judgfail
\end{align*}\noindent
To deal with this situation, our type system includes two ways to
combine typing contexts, $\ctcomma$ and $\ctstar$, each of which
affords different capabilities when reasoning about, e.g., subtyping
relationships. Informally, $\Gamma_1 \ctstar \Gamma_2$ means that the
capabilities of $\Gamma_1$ and $\Gamma_2$ are independent; when either
context includes an angelic type binding, our type system is free to
``choose'' a favorable assignment for the corresponding variable
independent of the other context.  To illustrate, consider the
following two subtyping judgments:
\begin{align*}
  &x{:}\urt{\Nat}{\vnu > 0} \ctstar y{:}\urt{\Nat}{\vnu > 1}
    \vdash \urt{\Nat}{\vnu = x - y} \subty \urt{\Nat}{\vnu > 0}
    \;\judgok
  &x{:}\urt{\Nat}{\vnu > 0} \ctcomma y{:}\urt{\Nat}{\vnu > 1}
    \vdash \urt{\Nat}{\vnu = x - y} \subty \urt{\Nat}{\vnu > 0} \;\judgfail
\end{align*}\noindent
Under the first context, we can freely choose both an arbitrary
positive value for $x$ \emph{and} a $2$ for $y$, from which the
subtyping relationship follows. Under the second context, however, the
possible values we can choose for $x$ and $y$ are \emph{Entangled}
conjunction, $\ctcomma$, meaning we are free to choose a value for one
but not the other. Intuitively, the hole in \ref{ex:EnCon} has this
typing context: once $x$ is fixed, the value of $y$ depends on it, and
vice-versa.  As a consequence, we are unable prove the following
typing judgement:
\begin{align*}
  &x{:}\urt{\Nat}{\vnu > 0} \ctcomma y{:}\urt{\Nat}{\vnu > 0} \vdash x - y : \urt{\Nat}{\vnu > 0} \;\judgfail
\end{align*}\noindent

The \ctcomma in the context
$\Gamma \cdot x{:}\tau_x \ctcomma y{:}\tau_y$ works like additive
logic conjunction in BI, and both $\Gamma \vdash x{:}\tau_x$ and
$\Gamma \vdash y{:}\tau_y$ are valid. Importantly, however, does not
guarantee that the modalities of $x$ and $y$ hold \emph{at the same
  time}. \footnote{Unlike a linear type system, which would
  \emph{consume} the capability to angelically assign $x$ when
  producing the capability to angelically assign $y$, our system
  instead \emph{entangles} these capabilities, allowing the body
  expression $e$ to use both $x$ and $y$ while still tracking their
  dependency. This formulation aligns with the intuition of bunched
  typing~\cite{BunchedType}. \BD{I don't think the reader is going to
    appreciate this difference unless they have a better intuition of
    the capabilities afforded by $\ctcomma$.}}  For example,
\begin{align*}
  &x{:}\urt{\Nat}{\vnu > 0} \ctstar (y{:}\urt{\Nat}{\vnu > 0} \ctcomma z{:}\urt{\Nat}{\vnu > 0}) \vdash x - y : \urt{\Nat}{\vnu > 0} \;\judgok
  \\&x{:}\urt{\Nat}{\vnu > 0} \ctstar (y{:}\urt{\Nat}{\vnu > 0}
  \ctcomma z{:}\urt{\Nat}{\vnu > 0}) \vdash x - (z + y) : \urt{\Nat}{\vnu > 0} \;\judgok
  \\&x{:}\urt{\Nat}{\vnu > 0} \ctstar (y{:}\urt{\Nat}{\vnu > 0} \ctcomma z{:}\urt{\Nat}{\vnu > 0}) \vdash y - z : \urt{\Nat}{\vnu > 0} \;\judgfail
\end{align*}\noindent
In these judgements, the context tracks variable dependency,
specifying that $x$ is independent from $y$ and $z$, while $y$ and $z$
are dependent on each other.  Specifying independence in this way is
sufficient to verify reachability for $x - y$ and $x - z$, since $x$
can be angelically chosen independently of $y$ or $z$, resp.  However,
type-checking $y - z$ against the type $\urt{\Nat}{\vnu > 0}$ fails
because $y$ and $z$ are entangled: no angelic choice can
simultaneously make their difference reach every positive number.

Notably, angelic and demonic refinements can also be entangled in a
way that impacts the capabilities afforded by a type context, as the
following pair of judgments illustrate:
\begin{align*}
  &x{:}\urt{\Nat}{\top} \ctcomma y{:}\ort{\Nat}{\top} \vdash x - y : \urt{\Nat}{\I{even}(\vnu)} \;\judgfail
  \\&x{:}\urt{\Nat}{\top} \ctstar y{:}\ort{\Nat}{\top} \vdash x - y : \urt{\Nat}{\I{even}(\vnu)} \;\judgok
\end{align*}\noindent
The first typing context stipulates that the value of $y$ is positive,
but that value is allowed to depend on the value $x$ --- the hole in
\ref{ex:EnCon} can also be typed under this typing
context. Accordingly, the expression \ref{ex:EnCon}$[x - y]$ cannot
reach every even number. Contrast this with a variant of this context
that instead assigns \emph{some} constant value $n$ to $y$:
\begin{align}
  &\zlet{x}{\Code{nat\_gen () }}{\zlet{y}{n}{\square}}
    \tag{C$_{n}$}\label{ex:DisjSafeCon} \\
\end{align}
Regardless of the particular value of $n$, for any given even number
$m$, choosing $x = m + n$ allows \ref{ex:DisjSafeCon}$[x - y]$ to
reach that number.

% \BD{The $\ctsubseteqA$ relation is doing some heavy lifting in our
%   type system, but at this point, it is not clear how that relation
%   interacts with additive and multiplicative conjunction. Maybe we can
%   provide some algebraic laws to help build up the reader's
%   intuitions?}

\paragraph{Function Types}
\BD{We need to (re-)emphasize here that this is the first refinement
  type system that does not place any restrictions on where angelic
  and demonic modalities can appear in function types. }

Entangled ($\ctsarr$) and independent
function types ($\ctwand$) are adjunctions of the corresponding type
context composition operators:
\begin{align*}
  \Gamma \vdash e_f : x{:}\tau_x \ctsarr \tau \quad &\text{ iff } \quad \Gamma \ctcomma x{:}\tau_x \vdash (e_f\;x) : \tau
  \\\Gamma \vdash e_f : x{:}\tau_x \ctwand \tau \quad &\text{ iff } \quad \Gamma \ctstar x{:}\tau_x \vdash (e_f\;x) : \tau
\end{align*}\noindent
Thus, ($\Code{-}_{\Code{nat}}$) can be given the following types:
\begin{align*}
  &\vdash (-) : x{:}\urt{\Nat}{\vnu > 0} \ctwand y{:}\urt{\Nat}{\vnu > 0} \ctwand \urt{\Nat}{\vnu > 0}
  \\&\vdash (-) : x{:}\urt{\Nat}{\vnu > 0} \ctsarr y{:}\urt{\Nat}{\vnu > 0} \ctwand \urt{\Nat}{\vnu > 0}
\end{align*}\noindent
Notably, it cannot have the following type:
\begin{align*}
  &\vdash (-) : x{:}\urt{\Nat}{\vnu > 0} \ctwand y{:}\urt{\Nat}{\vnu > 0} \ctsarr \urt{\Nat}{\vnu > 0}
\end{align*}\noindent
Doing so would requires proving the following type judgement:
\begin{align*}
  &\emptyset \ctstar x{:}\urt{\Nat}{\vnu > 0} \ctcomma y{:}\urt{\Nat}{\vnu > 0} \vdash x - y : \urt{\Nat}{\vnu > 0} \;\judgfail
\end{align*}\noindent
As we saw earlier, a type context in which $x$ and $y$ are entangled
restricts the angelic choices that can be made, preventing us from
establishing that $x - y$) can produce every positive number.

Higher-order functions introduce additional complexity. Consider the
following example:
\begin{align*}
  &x{:}\urt{\Nat}{\vnu > 0} \ctstar y{:}\urt{\Nat}{\vnu > 0} \vdash \zzlam{f}{x - (f\; (f\;y))} :
  \\&\quad  \underbrace{f{:}(\underbrace{y{:}\urt{\Nat}{\vnu > 0} \ctsarr}_\texttt{$y$ can depend on context during application} \urt{\Nat}{\vnu > 0}) \ctwand}_\texttt{$f$ is independent from $x$ and $y$} \urt{\Nat}{\vnu > 0}
\end{align*}\noindent
To make $x - (f\;(f\;y))$ reach every positive number, $x$ and
$f\;(f\;y)$ must be independent.  To ensure this, one might expect
that all function arrows in the result type would need to be
$\ctwand$, but this would be incorrect. If $f$ had the fully
independent type
$y{:}\urt{\Nat}{\vnu > 0} \ctwand \urt{\Nat}{\vnu > 0}$, then each
application of $f$ would require its argument to be independent from
$f$'s closure. While the first application $f\;y$ satisfies this, the
second application $f\;(f\;y)$ fails because the intermediate result
$f\;y$ depends on $f$.  While $\ctwand$ guarantees that $f$'s argument
is not entagled with its closure, it makes no guarantees that $f$'s
result is also independent.  Thus, while $(f\; y)$ is independent of
$y$, it may still depend on $f$'s closure, making it an invalid
argument in the outer application when equipped with this signature.

If $f$ is instead given the entangled function type
$y{:}\urt{\Nat}{\vnu > 0} \ctsarr \urt{\Nat}{\vnu > 0}$, its argument
can now depend on $f$'s closure, accommodating the nested
application. The function itself is then typed with $\ctwand$,
ensuring $f$'s body is not entagled with either $x$ or $y$ in the
outer context. By adjunction, this yields:
\begin{align*}
  &x{:}\urt{\Nat}{\vnu > 0} \ctstar
  \\&\quad (y{:}\urt{\Nat}{\vnu > 0} \ctstar
  f{:}(y{:}\urt{\Nat}{\vnu > 0} \ctsarr \urt{\Nat}{\vnu > 0}))
  \vdash x - (f\; (f\;y)) : \urt{\Nat}{\vnu > 0}
\end{align*}\noindent
Since $x$ is independent from both $y$ and $f$, it is also independent
from $f\;(f\;y)$, thus guaranteeing that $x - (f\;(f\;y))$ can reach
every positive number.

\paragraph{From Reachability to Safety}
As discussed above, the sort of safety that can be directly derived
from reachability assumptions in the typing context is quite limited
--- ascribing a demonic refinement to a term constrains \emph{all} of
its possible executions, while an angelic type binding only makes
assumptions about a \emph{subset} of the contexts that term may be run
in. It is possible, however, to derive meaningful safety properties
about the executions that an angelic refinement guarantees exist. If
we have assumed that $x$ must reach every even number, for example, we
know that the term $x \Code{+ 1}$ will always safely evaluate to an
odd number \emph{under those environments}. In order to soundly
establish upper bounds on program executions that are guaranteed to
exist, our type system includes a \emph{context sum} type constructor
$\ctoplus$ that explicitly identifies partitions of the values a term
must reach.

The typing judgement $\Gamma \vdash e : \tau_1 \ctoplus \tau_2$ states
that $e$ can be typed with $\tau_1$ and $\tau_2$ in some partitioning
of $\Gamma$, and that all environments consistent with $\Gamma$ are
accounted for by $\tau_1$ and $\tau_2$ --- no possibility is lost.
For example:
\begin{align*}
  &x{:}\urt{\Nat}{even(x)} \vdash x + 1 : \ort{\Nat}{odd(\vnu)} \ctoplus
    \ort{\Nat}{even(x)} \;\judgok \\
  &x{:}\urt{\Nat}{\top} \vdash \zif{x == 0}{3}{5} : \ort{\Nat}{\vnu < 4}\ctoplus \ort{\Nat}{\vnu > 4} \;\judgok
\end{align*}\noindent
In both judgements, the result type uses $\ctoplus$ to explicitly
demarcate two subsets of the term's executions: the first example uses
$\ctoplus$ to distinguish between the executions that must exist due
to the angelic refinement of $x$ in the typing context that those that
may exist, while the second uses $\ctoplus$ to distinguish between the
two possible control flows through the program.

Importantly, while
$\ort{\Nat}{\vnu < 4} \ctoplus \ort{\Nat}{\vnu > 4}$ is a subtype of
$\ort{\Nat}{\vnu < 4 \lor \vnu > 4}$ (both types can only be ascribed
to a term that never evaluates to $4$), the reverse is not true: the
disjunctive refinement type loses the information that a term has two
distinct control flow paths, one whose result is always greater than
$4$ and another whose result is always less than $4$. \BD{We might
  want to nod to outcome logic here.} This ability also allows us to
focus on the safety properties of a particular program path, assigning
the remainder the $\top$ type:
\begin{align*}
  &x{:}\urt{\Nat}{\top} \vdash \zif{x > 0}{5}{3} : \ort{\Nat}{\vnu > 4} \ctoplus \ort{\Nat}{\top} \quad \judgok
  \\&x{:}\urt{\Nat}{\top} \vdash \zif{x > 0}{5}{3} : \ort{\Nat}{\vnu < 4} \ctoplus \ort{\Nat}{\top} \quad \judgok
\end{align*}\noindent
The use of $\tau \ctoplus \ort{b}{\top}$ expresses an \emph{affine}
style of control flow analysis when deriving information from angelic
facts: like affine logic, which allows resources to be used at most
once (and hence discarded), this style allows the type system to
discard uninteresting branches by weakening them to $\ort{b}{\top}$,
imposing a precise safety constraint on the branches that a typing
context guarantees exist while placing no constraint on the others.

% Standard refinement types interpret type
% contexts in a demonic style, requiring \emph{all} control-flow within
% the program to be consistent with the safety property.  In the
% presence of angelic contexts that can reach multiple control flows,
% additional handling is required.  Consider the following judgements:
% \begin{align*}
%   &x{:}\urt{\Nat}{\top} \vdash \zif{x == 0}{3}{5} : \ort{\Nat}{\vnu < 4} \quad \text{(reach true branch)}
%   \\&x{:}\urt{\Nat}{\top} \vdash \zif{x == 0}{3}{5} : \ort{\Nat}{\vnu > 4} \quad \text{(reach false branch)}
% \end{align*}\noindent
% These judgements informally state that ``there exists an input natural
% number $x$ such that we may reach some number less than (greater than,
% resp.) $4$.''  They mix safety and reachability reasoning, encoding a
% reachability property using imprecise overapproximated constraints
% (i.e., $\ort{\Nat}{\vnu < 4}$ and $\ort{\Nat}{\vnu > 4}$) instead of
% precise underapproximated constraints (i.e., $\urt{\Nat}{\vnu = 3}$
% and $\urt{\Nat}{\vnu = 5}$), since the desired safety property does not need
% to be precise about the exact result values.

% However, na\"{i}vely combining these imprecise reachability properties
% with safety reasoning may lead to semantic conflicts.  A purely safety-focussed
% interpretation of the first judgement would require that
% \emph{every} value produced by $\zif{x == 0}{3}{5}$ under the
% angelic context $x{:}\urt{\Nat}{\top}$ satisfies $\vnu < 4$.  This is
% incorrect: since $x$ can angelically take any natural number,
% including $1$, the expression can evaluate to $5$, violating $\vnu < 4$.
% This conflict also breaks the standard semantic characterization
% of intersection types found in refinement type systems that are used to model
% conditionals:
% \begin{align*}
%   &\Gamma \vdash e : \ort{b}{\phi_1} \ctinter \ort{b}{\phi_2} \text{ iff } \Gamma \vdash e : \ort{b}{\phi_1 \land \phi_2}
% \end{align*}\noindent
% since no term can have type $\ort{\Nat}{\vnu < 4 \land \vnu > 4}$.
% The root cause is that an angelic context cannot be treated as simple
% existential quantification (e.g., ``there exists an $x$ that reaches
% the true branch'') that discards the possibility that $x$ can also
% reach the false branch.  In order to preserve all possibilities while
% still enabling an existential-like interpretation of angelic contexts,
% we introduce a \emph{context sum} type constructor $\ctoplus$ that
% explicitly tracks  expected control flows.

\section{Context Type System}
\label{sec:context-typing}
% In this section, we introduce the core language and the context type system for it.

% \subsection{Core Language}

\begin{figure}[h]
  {\small
    \begin{alignat*}{2}
      \text{\textbf{Variables }}& \quad &\quad& x, y, z, \vnu, ...
      \\\text{\textbf{Data Constructors }}& \quad &d ::= \quad & () ~|~ \Code{true} ~|~ \Code{false} ~|~ \Code{O} ~|~ \Code{S} ~|~ \Code{Cons} ~|~ \Code{Nil} ~|~ ...
      \\\text{\textbf{Operations}}& \quad & \primop ::= \quad & d ~|~ {+} ~|~ {-} ~|~ {==} ~|~ {<} ~|~ {\leq} ~|~ ...
      \\ \text{\textbf{Constants }}& \quad &c ::= \quad & () ~|~ \mathbb{B} ~|~ \mathbb{Z} ~|~ d\ \overline{c}\ ...
      \\ \text{\textbf{Values}}& \quad  & v ::= \quad & c ~|~ x ~|~ \zlam{x}{s}{e} ~|~ \zfix{f}{s}{x}{s}{e}
      \\ \text{\textbf{Expressions}}& \quad & e ::=\quad & v ~|~ \zlet{x{:}s}{e}{e} ~|~ \primop\ \overline{v} ~|~  v\;v ~|~ \match{v} \overline{d\ \overline{y} \to e}
      \\\text{\textbf{Base Types}}& \quad & b  ::= \quad &  \Code{unit} ~|~ \Code{bool} ~|~ \Code{nat} ~|~ \List[b] ~|~ ...
      \\\text{\textbf{Basic Types}}& \quad & s  ::= \quad &  b ~|~ s\sarr s
      \\\text{\textbf{Basic Type Context}}& \quad & \Sigma ::= \quad &  \emptyset ~|~ x{:}s ~|~ \Sigma, \Sigma
    \end{alignat*}}
  \caption{Syntax of \langname{} \BD{Why do we include both constants
      and data constructors?}}
  \label{fig:core-lang-syntax}
  \end{figure}

  % \paragraph{Terms, basic types, and basic type contexts}
  We formalize our proposed context type system for the core language
  (\langname{}) shown in \autoref{fig:core-lang-syntax}.  This
  language is a standard call-by-value lambda calculus with inductive
  datatypes, pattern matching (the following uses conditional
  expressions as syntactic sugar for \Code{match}), and recursive
  functions. \BD{Are we assuming that built-in operations are
    deterministic? Do we even need these?}  Expressions are written in
  monadic normal form (MNF)~\cite{JO94}, a variant of A-Normal Form
  (ANF)~\cite{FSDF93} that allows nested let-bindings.  The language's
  basic types (i.e., $s$) consist of base types ($b$) and function
  types (i.e., $s\sarr s$). Both \textbf{let}-expressions and function
  abstractions explicitly annotate bound variables with a basic type
  (e.g., $\zlam{x}{s}{e}$); we omit these annotations for brevity,
  when appropriate.  A basic type context $\Sigma$ maps variables to
  their basic types; we assume that all bound variables in $\Sigma$
  are distinct.

\begin{figure}[h]
  { \small
    \begin{alignat*}{2}
      \text{\textbf{Qualifiers}}& \quad & \phi ::= \quad & c ~|~ x ~|~ \primop\;\overline{v} ~|~ \bot ~|~ \top  ~|~ \neg \phi ~|~ \phi \land \phi ~|~ \phi \lor \phi ~|~  \phi \impl \phi ~|~ \forall x{:}b. \phi ~|~ \exists x{:}b. \phi
      \\ \text{\textbf{Refinement Types}}& \quad & \tau ::= \quad & \ort{b}{\phi} ~|~ \urt{b}{\phi} ~|~ \tau \ctinter \tau ~|~ \tau \ctunion \tau ~|~ \tau \ctoplus \tau ~|~ x{:}\tau \ctsarr \tau ~|~ x{:}\tau \ctwand \tau
      % ~|~ \ctpers \tau
      \\ \text{\textbf{Type Contexts}}& \quad & \Gamma ::= \quad & \emptyset ~|~ x{:}\tau ~|~ \Gamma \ctstar \Gamma ~|~ \Gamma \ctcomma \Gamma ~|~ \Gamma \ctoplus \Gamma
      % \\ \text{\textbf{Contexts With Holes}}& \quad & \Delta ::= \quad & \Box ~|~ \Gamma ~|~ \Delta \ctstar \Delta ~|~ \Delta \ctcomma \Delta ~|~ \Delta \ctoplus \Delta
    \end{alignat*}}
  %% \vspace*{-.1in}
  %% {\small
  %% \begin{align*}
  %%   \text{\textbf{Type erasure}} \quad & \eraserf{\ort{b}{\phi}} \doteq b \quad \eraserf{\urt{b}{\phi}} \doteq b \quad \eraserf{\tau_1 \ctinter \tau_2} \doteq \eraserf{\tau_1} \quad \eraserf{\tau_1 \ctunion \tau_2} \doteq \eraserf{\tau_1}
  %%   \\&\eraserf{\tau_1 \ctoplus \tau_2} \doteq \eraserf{\tau_1} \quad  \eraserf{x{:}\tau_x \ctsarr \tau} \doteq \eraserf{\tau_x} \sarr \eraserf{\tau} \quad
  %%   \\& \eraserf{\emptyset} \doteq \emptyset \quad \eraserf{x{:}\tau} \doteq x{:}\eraserf{\tau} \quad \eraserf{\Gamma_1 \ctcomma \Gamma_2} \doteq \eraserf{\Gamma_1}, \eraserf{\Gamma_2} \quad \eraserf{\Gamma_1 \ctstar \Gamma_2} \doteq \eraserf{\Gamma_1}, \eraserf{\Gamma_2}  \quad \eraserf{\Gamma_1 \ctoplus \Gamma_2} \doteq \eraserf{\Gamma_1}
    % \\\text{\textbf{Type lifting}} \quad & \liftrf{b} \doteq \ort{b}{\phi} \quad \liftrf{s_x \sarr s} \doteq x{:}\liftrf{s_x} \ctsarr \liftrf{s}
    % \\&\liftrf{\emptyset} \doteq \emptyset \quad \liftrf{x{:}s} \doteq x{:}\liftrf{s} \quad \liftrf{\Gamma_1 , \Gamma_2} \doteq \liftrf{\Gamma_1} \ctcomma \liftrf{\Gamma_2}
%  \end{align*}}
  \caption{Context Type Syntax}
  \label{fig:context-types-syntax}
\end{figure}

The syntax of context types is shown in
\autoref{fig:context-types-syntax}. As in other refinement type
systems, \emph{context types} refine basic types with qualifiers
$\phi$; but each refinement is ascribed either a demonic and angelic
modality ($\otparenth{...}$ and $\utparenth{...}$, resp.). Context
types additionally include bunched implication (e.g., $\ctsarr$,
$\ctwand$), disjoint sums ($\ctoplus$), and standard type intersection
($\ctinter$) operators. Our type system shares the same intuition
underlying bunched implication and refinement types, generally;
however, its technical development differs significantly from both due
to its support for angelic contexts and mixed demonic/angelic
reasoning.

\begin{figure}[h!]
  {\small
  {\normalsize \begin{flalign*}
    &\text{\textbf{Well-formedness rules}} && \fbox{$\Sigma \wfvdash \tau \quad \Sigma \wfvdash \Gamma$}
  \end{flalign*}}
% \mprooftr{28mm}
% {$\Gamma \basicvdash \tau$ \quad $\wfvdash \Gamma$}
% {WfType}
% {$\Gamma \wfvdash \tau$}
% \quad
\mprooftr{35mm}
{$\Sigma \wfvdash \Gamma_1$ \quad $\Sigma, \eraserf{\Gamma_1} \wfvdash \Gamma_2$}
{WfComma}
{$\Sigma \wfvdash \Gamma_1 \ctcomma \Gamma_2$}
\quad
\mprooftr{35mm}
{$\Sigma \wfvdash \Gamma_1$ \quad $\Sigma \wfvdash \Gamma_2$}
{WfStar}
{$\Sigma \wfvdash \Gamma_1 \ctstar \Gamma_2$}
% \mprooftr{53mm}
% {$\forall r. r \models \denot{\Gamma} \impl \forall e. r \models \denot{\tau_1}\;e \impl r \models \denot{\tau_2}\;e$}
% {Sub}
% {$\Gamma \vdash \tau_1 <: \tau_2$}
% \quad
% \mprooftr{50mm}
% {$\forall r. \denot{\Gamma_1}\;r \impl \exists r'. \restrictTo{r}{X} \sqsubseteq r' \land \denot{\Gamma_2}\;r'$ }
% {CtxSub}
% {$\Gamma_1 \ctsubseteq{X} \Gamma_2$}
	  }
  \caption{Selected auxiliary well-formedness relations\footnote{Well-formedness rules for types
      and contexts check that qualifiers are well-typed with respect to the basic
      type structure.  The $\eraserf{.}$ operator converts context types $\tau$ to basic types $s$
      by stripping away all refinement information.  Details are provided in the supplemental material.}}
  \label{fig:auxiliary-rules}
  \vspace*{-.1in}
\end{figure}

\paragraph{Well-formedness} Similar to other refinement type systems, our type system depends on the well-formedness relations shown in
\autoref{fig:auxiliary-rules}. The well-formedness of types ($\Sigma
\wfvdash \tau$) and type contexts ($\Sigma \wfvdash \Gamma$) mostly
requires type qualifiers to be well-typed in the basic type system and
binding variables in type contexts to be unique. We would like to
highlight that rule \textsc{WfComma} allows the second context to
refer to variables in the first context, while rule \textsc{WfStar}
requires that $\Gamma_1$ and $\Gamma_2$ are completely disjointed. For
example, within a type context $\Gamma_0 \ctcomma (\Gamma_1 \ctstar
\Gamma_2)$, context $\Gamma_0$ is closed, and $\Gamma_1$ and
$\Gamma_2$ can only refer to variables in $\Gamma_0$.


\begin{figure}[h!]
  {\small
  {\normalsize \begin{flalign*}
    &\text{\textbf{Subtyping rules}} && \fbox{$\Gamma \vdash \tau \subty \tau$ \quad $\Gamma \ctsubseteq{X} \Gamma$ }
    \\&\text{\textbf{Typing rules}} && \fbox{$\Gamma \vdash e : \tau$}
  \end{flalign*}}
  \mprooftr{40mm}
  {$\Gamma \vdash e_1 : \tau_1 $ \quad $\Gamma \ctcomma x{:}\tau_1 \vdash e_2 : \tau_2$}
  {T-Let}
  {$\Gamma \vdash \zlet{x}{e_1}{e_2} : \tau_2$}
  \quad
  \mprooftr{40mm}
  {$\Gamma_1 \vdash e_1 : \tau_1 $ \quad $\Gamma_2 \ctstar x{:}\tau_1 \vdash e_2 : \tau_2$}
  {T-LetD}
  {$\Gamma_1 \ctstar \Gamma_2 \vdash \zlet{x}{e_1}{e_2} : \tau_2$}
  \\[.8em]
  \mprooftr{42mm}
  {$\eraserf{\tau_x} = s_x$ \quad $\Gamma \ctcomma x{:}\tau_x \vdash e : \tau$}
  {T-Lam}
  {$\Gamma \vdash \zlam{x}{s_x}{e} : x{:}\tau_x \ctsarr \tau$}
  \quad
  \mprooftr{42mm}
  {$\eraserf{\tau_x} = s_x$ \quad  $\Gamma \ctstar x{:}\tau_x \vdash e : \tau$}
  {T-LamD}
  {$\Gamma \vdash \zlam{x}{s_x}{e} : x{:}\tau_x \ctwand \tau$}
  \\[.8em]
  \mprooftr{40mm}
  {$\Gamma \vdash v_1 : x{:}\tau_x \ctsarr \tau$ \quad
  $\Gamma \vdash v_2 : \tau_x$ }
  {T-AppFun}
  {$\Gamma \vdash v_1\;v_2 : \tau[x \mapsto v_2]$}
  \quad
  \mprooftr{46mm}
  {$\Gamma_1 \vdash v_1 : x{:}\tau_x \ctwand \tau$ \quad
  $\Gamma_2 \vdash v_2 : \tau_x$  }
  {T-AppFunD}
  {$\Gamma_1 \ctstar \Gamma_2 \vdash v_1\;v_2 : \tau[x \mapsto v_2]$}
  \\[.8em]
  \mprooftr{30mm}
  {$\Gamma \vdash e : \tau_1$ \quad $\Gamma \vdash \tau_1 \subty \tau_2$}
  {T-Sub}
  {$\Gamma \vdash e : \tau_2$}
  \quad
  \mprooftr{40mm}
  {$\Gamma_2 \vdash e : \tau$ \quad $\Gamma_1 \ctsubseteq{\Code{FV}(e) \cup \Code{FV}(\tau)} \Gamma_2$}
  {T-CtxSub}
  {$\Gamma_1 \vdash e : \tau$}
  \quad
  \mprooftr{14mm}
  {$ $}
  {T-Var}
  {$x{:}\tau \vdash x : \tau$}
  \\[.8em]
  \mprooftr{43mm}
  {$ \emptyset \basicvdash c : b $}
  {T-Const}
  {$\emptyset \vdash c : \ort{b}{\vnu = c} \ctinter \urt{b}{\vnu = c}$}
  \quad
  \mprooftr{54mm}
  {$\Gamma \ctcomma x{:}\tau_x \ctcomma f{:}(x{:}\ort{b}{\phi \land \vnu \prec x} \ctsarr \tau) \vdash e : \tau$}
  {T-Fix}
  {$\Gamma \vdash \zzfix{f}{x}{e} : x{:}\ort{b}{\phi} \ctsarr \tau$}
  \\[.8em]
  \mprooftr{34mm}
  {$\Code{Ty}(\primop) = \overline{x_i{:}\tau_i}\ctsarr \tau_y$ \quad
    $\forall i. \Gamma \vdash v_i : \tau_i$ }
  {T-AppOp}
  {$\Gamma \vdash \primop\;\overline{v_i} : \tau[\overline{x_i \mapsto v_i}]$}
  \quad
  \mprooftr{64mm}
  {$\forall i. \Gamma_i \vdash d_i(\overline{y}) : \ort{b}{\vnu = v} \interty \urt{b}{\vnu = v}$ and $\Gamma_i \vdash e_i : \tau_i$ }
  {T-Match}
  {$\ctbigoplus \overline{\Gamma_i} \vdash \match{v} \overline{d\ \overline{y} \to e}  : \ctbigoplus \overline{\tau_i}$}
  }
  \caption{Declarative typing rules for the core language.}
  \label{fig:declarative-typing-rules}
  \vspace*{-.1in}
\end{figure}

\paragraph{Variables and Constants} The typing rule for variables is standard;
a constant is allowed to have the most precise (singleton) type, which
has the same interpretation in both a demonic as well as angelic
setting.  In other words, the singleton type $\ort{b}{\vnu = c}
\ctinter \urt{b}{\vnu = c}$ coincides in both modalities because
there is only one possible value, so the safety and reachability
constraints are identical.


\paragraph{Bunched Typing} The typing rules are shown in \autoref{fig:declarative-typing-rules},
where rules for let-expressions, lambda functions, and applications
have dependent and disjointed versions, respectively.  Unlike standard
bunched typing, the dependent version does not require an explicit
additive conjunction in the type context, as shown below:
\begin{align*}
\mprooftr{40mm}
  {$\Gamma_1 \vdash e_1 : \tau_1 $ \quad $\Gamma_2 \ctcomma x{:}\tau_1 \vdash e_2 : \tau_2$}
  {T-Bunched-Let}
  {$\Gamma_1 \ctcomma \Gamma_2 \vdash \zlet{x}{e_1}{e_2} : \tau_2$}
\qquad
\mprooftr{26mm}
  {$\Gamma_1 \ctcomma \Gamma_2 \ctcomma \Gamma_2 \vdash e : \tau$}
  {T-Bunched-Contraction}
  {$\Gamma_1 \ctcomma \Gamma_2 \vdash e : \tau$}
\end{align*}\noindent
Unlike standard bunched typing, the dependent version does not require
an explicit additive conjunction in the type context. The
well-formedness condition requiring unique variable bindings means
that the substructural contraction rule
\textsc{T-Bunched-Contraction}, which duplicates bindings, is
incompatible with our system. Instead, our \textsc{T-Let} rule
directly entangles the newly added dependent context with the existing
context, without duplicating or consuming bindings.


\begin{example}[Let-expression] Consider the following example of typing for a let-expression:
\begin{align*}
    x{:}\ort{\Nat}{\vnu < 10} \vdash \zlet{y}{x - 3}{x - y} : \ort{\Nat}{\vnu \leq 3}
\end{align*}\noindent
Note that binding term $x - 3$ and body expression $x - y$ share the same variable $x$. If we apply standard bunched typing rule \textsc{T-Bunched-Let}, we require the following type context:
\begin{align*}
  \\&\underbrace{x{:}\ort{\Nat}{\vnu < 10}}_{\texttt{for $x - 3$}} \ctcomma \underbrace{x{:}\ort{\Nat}{\vnu < 10}}_{\texttt{for $x - y$}} \vdash \zlet{y}{x - 3}{x - y} : \ort{\Nat}{\vnu \leq 3}
\end{align*}\noindent
which requires using rule \textsc{T-Bunched-Contraction} to ``duplicate'' binding $x$. However, our \textsc{T-Let} rule optimizes this ``duplicate-then-consume'' pattern to yield the following derivation:
\begin{align*}
  &x{:}\ort{\Nat}{\vnu < 10} \vdash x - 3 : \ort{\Nat}{\vnu < 7}
  \\&\underline{x{:}\ort{\Nat}{\vnu < 10} \ctcomma y{:}\ort{\Nat}{\vnu < 7} \vdash x - y : \ort{\Nat}{\vnu \leq 3}}
  \\&x{:}\ort{\Nat}{\vnu < 10} \vdash \zlet{y}{x - 3}{x - y} : \ort{\Nat}{\vnu \leq 3}
\end{align*}\noindent
\end{example}

\paragraph{Subtyping}  Subsumption rules for both types (\textsc{T-Sub}) and
type contexts (\textsc{T-CtxSub}) use a \emph{semantic
  subtyping} relation. The type subtyping relation $\Gamma
\vdash\tau_1 \subty \tau_2$ is standard, i.e., inhabitants of the
denotation of $\tau_1$ are also inhabitants of the denotation of
$\tau_2$ under the denotation of context $\Gamma$. The context
subtyping relation $\Gamma_1 \ctsubseteq{X} \Gamma_2$ follows the
intuition in previous examples, where we restrict comparison to a set
of variables $X$.  For example,
\begin{align*}
  x{:}\ort{\Nat}{\vnu > 0} \ctcomma y{:}\urt{\Nat}{\top}
  \ctsubseteq{\{y\}} y{:}\urt{\Nat}{\top} \ctcomma z{:}\urt{\Nat}{\top}
\end{align*}
Rule \textsc{T-CtxSub} only requires the context-subtyping relation
to hold over free variables of the term and type. Further details are
given in~\autoref{sec:meta}.

\begin{example}[Substructural rules] We can derive the following substructural rules from \textsc{T-CtxSub}:
  \begin{align*}
    \mprooftr{25mm}
      {$\Gamma_2 \vdash e : \tau$}
      {T-Frame}
      {$\Gamma_1 \ctstar \Gamma_2 \vdash e : \tau$}
    \quad
    \mprooftr{26mm}
      {$\Gamma_2 \vdash e : \tau$}
      {T-Weakening}
      {$\Gamma_1 \ctcomma \Gamma_2 \vdash e : \tau$}
    \quad
    \mprooftr{25mm}
      {$\Gamma_2 \ctstar \Gamma_1 \vdash e : \tau$}
      {T-Exchange}
      {$\Gamma_1 \ctstar \Gamma_2 \vdash e : \tau$}
    \end{align*}\noindent
The variable-set constraint in context subtyping relation is important.
For example, framing rule \textsc{T-Frame} actually requires that $\Gamma_1 \ctstar \Gamma_2 \ctsubseteq{\Code{FV}(e) \cup \Code{FV}(\tau)} \Gamma_2$ holds, which is silently forced by $\Gamma_2 \vdash e : \tau$, requiring that the domain of $\Gamma_2$ is a subset of $\Code{FV}(e) \cup \Code{FV}(\tau)$.
\end{example}

\paragraph{Fixpoint Functions} As in standard refinement type systems,
we require all well-typed terms to be total. Thus, rule \textsc{T-Fix}
forces its first argument to have only base type and to always
decrease according to some well-founded relation $\prec$. This
decreasing constraint means the first parameter type cannot have
angelic modality; otherwise, the argument of recursive call
$v_{\Code{arg}}$ is required to have the following type:
\begin{align*}
  x{:}\urt{b}{\phi} \vdash v_{\Code{arg}} : \urt{b}{\phi} \interty
  \ort{b}{\vnu \prec x}
\end{align*}\noindent
where type $\urt{b}{\phi}$ requires $v_{\Code{arg}}$ to reach all
values that $x$ can reach; however, type $\ort{b}{\vnu \prec x}$
disallows $v_{\Code{arg}}$ from reaching the value of $x$. Such an
argument does not exist, which means recursive calls can never happen
when the parameter type has an angelic modality. Moreover, the
parameter type for a recursive function cannot be a disjointed arrow
type, i.e.,
\begin{align*}
  \Gamma \ctcomma x{:}\tau_x \ctcomma f{:}(x{:}\ort{b}{\phi \land \vnu \prec x} \ctwand \tau) \vdash e : \tau \quad \judgfail
\end{align*}\noindent
The parameter type of $f$ in a recursive call must reference $x$ via
the well-foundedness constraint $\nu \prec x$, which establishes a
dependency between the argument of the recursive call and the original
parameter $x$.  This dependency makes the disjointness requirement of
$\ctwand$ impossible to satisfy, since $\ctwand$ requires the argument
context to be independent of the closure context, whereas well-founded
induction inherently requires the opposite. Although \textsc{T-Fix}
therefore restricts the function type of recursive calls to use $\ctsarr$,
the subsumption rule \textsc{T-Sub} still allows a recursive function
to be typed against arbitrary function types after the fact, provided
the subtyping relation holds.

\begin{example}[Fixpoint Function] The following fixpoint function defines natural number subtraction discussed in \autoref{sec:intro}.
\begin{align*}
  \vdash &\zzfix{f}{x}{\zzlam{y}{\zif{x = 0}{0}{\zif{y = 0}{x}{f\;(x - 1)\;(y - 1)}}}}
  \\& : x{:}\ort{\Nat}{\top} \ctsarr y{:}\ort{\Nat}{\top} \ctsarr \ort{\Nat}{\vnu = x -_{\Nat} y}
\end{align*}\noindent
With a stronger function type that enables type checking with rule \textsc{T-Fix}, we can convert it into the type shown in \autoref{sec:intro} via  \textsc{T-Sub}.
\begin{align*}
  \vdash ... :  x{:}\urt{\Nat}{\vnu > 0} \ctwand y{:}\urt{\Nat}{\vnu > 0} \ctwand \urt{\Nat}{\vnu > 0}
\end{align*}\noindent
\end{example}

\paragraph{Datatypes and pattern matching} The last two rules
in \autoref{fig:declarative-typing-rules} type data constructors and
primitive operators. Rule \textsc{T-AppOp} is similar to the
application rule but relies on an auxiliary function $\Code{Ty}$ to
assign types to primitive operators and datatype constructors in the
core language. We require these primitive types to be precise as in
\textsc{T-Const}; for example:
\begin{align*}
  \Code{Ty}(\Code{O}) &= \ort{\Nat}{\vnu = 0} \interty \urt{\Nat}{\vnu = 0}
  \\\Code{Ty}(\Code{S}) &= n{:}\ort{\Nat}{\top} \ctsarr \ort{\Nat}{\vnu = n+1} \interty \urt{\Nat}{\vnu = n + 1}
  \\\Code{Ty}(+) &= x{:}\ort{\Nat}{\top} \ctsarr y{:}\ort{\Nat}{\top} \ctsarr \ort{\Nat}{\vnu = x+y} \interty \urt{\Nat}{\vnu = x + y}
\end{align*}\noindent
With these precise primitive types, rule \textsc{T-AppOp} can forward-infer constraints of operator application results. As discussed previously, rule \textsc{T-Match} requires type context to be divided into a sum of contexts that align with \emph{reachable} control flows, and then the sum of corresponding result types becomes the result type of the whole pattern-matching expression. The reachability condition is guaranteed by $\Gamma_i$ making datatype $d_i(\overline{y})$ have a precise type equal to the matched value $v$, which requires the types of variables $\overline{y}$ to have correct precise types for further type-checking of control-flow branches.

\begin{example}[Pattern matching expression] We revisit our previous example where we unfold the conditional expression into a pattern-matching expression.
\begin{align*}
  x{:}\urt{\Nat}{\top} \vdash \zmatchNat{x}{3}{y}{5} : \ort{\Nat}{\vnu < 4} \ctoplus \ort{\Nat}{\top}
\end{align*}
where we want to show that the angelic context $x{:}\urt{\Nat}{\top}$ can be split so that
the true branch is reachable and types with $\ort{\Nat}{\vnu < 4}$.  We first divide the
type context via subsumption rule \textsc{T-CtxSub}.
{\small\begin{align*}
  &x{:}\urt{\Nat}{\top} \ctsubseteq{\{x\}}
  \\&\quad \underbrace{x{:}\ort{\Nat}{\vnu = 0}\interty \urt{\Nat}{\vnu = 0}}_{\Gamma_{\Code{trueBranch}}} \ctoplus \underbrace{x{:}\ort{\Nat}{\vnu > 0}\interty \urt{\Nat}{\vnu > 0} \ctcomma y{:}\ort{\Nat}{\vnu = x - 1}}_{\Gamma_{\Code{falseBranch}}}
\end{align*}}\noindent
which exactly aligns the assumption about context reachability required by rule \textsc{T-Match}:
\begin{align*}
  &\Gamma_{\Code{trueBranch}} \vdash 0 : \ort{\Nat}{\vnu = 0} \interty \urt{\Nat}{\vnu = 0}
  \\&\Gamma_{\Code{falseBranch}} \vdash \Code{S}\;y : \ort{\Nat}{\vnu = x} \interty \urt{\Nat}{\vnu = x}
\end{align*}\noindent
We then have the following derivation:
\begin{align*}
  &\Gamma_{\Code{trueBranch}} \vdash 3 : \ort{\Nat}{\vnu < 4}
  \\&\underline{\Gamma_{\Code{falseBranch}} \vdash 5 : \ort{\Nat}{\top}}
  \\&\Gamma_{\Code{trueBranch}} \ctoplus \Gamma_{\Code{falseBranch}} \vdash  \zmatchNat{x}{3}{y}{5} : \ort{\Nat}{\vnu < 4} \ctoplus \ort{\Nat}{\top}
\end{align*}\noindent
On the other hand, the following derivation fails:
\begin{align*}
  x{:}\urt{\Nat}{\vnu > 0} \vdash \zmatchNat{x}{3}{y}{5} : \ort{\Nat}{\vnu < 4} \ctoplus \ort{\Nat}{\top} \judgfail
\end{align*}\noindent
Because context $x{:}\urt{\Nat}{\vnu > 0}$ admits an environment $\Envir(e) \doteq \zlet{x}{\Code{positive\_gen}()}{e}$ under which execution always follows the false branch, it cannot be decomposed into a sum of two contexts that separately witness the true and false branches.
Alternatively, we also enable standard refinement type judgement with unreachable control flows:
\begin{align*}
  &x{:}\ort{\Nat}{\vnu > 0} \ctcomma y{:}\ort{\Nat}{\vnu = x - 1}\vdash 5 : \ort{\Nat}{\vnu > 4}
  \\&x{:}\ort{\Nat}{\vnu > 0} \ctcomma y{:}\ort{\Nat}{\vnu = x - 1}\vdash \Code{S}\;y : \ort{\Nat}{\vnu = x} \interty \urt{\Nat}{\vnu = x}
  \\&\underline{x{:}\ort{\Nat}{\vnu > 0} \ctsubseteq{\{x\}} x{:}\ort{\Nat}{\vnu > 0} \ctcomma y{:}\ort{\Nat}{\vnu = x - 1}}
  \\&x{:}\ort{\Nat}{\vnu > 0} \vdash \zmatchNat{x}{3}{y}{5} : \ort{\Nat}{\vnu > 4} \quad \judgok
\end{align*}\noindent
where we adapt rule \textsc{T-Match} with a sum of a single type.
\end{example}
